\chapter{综合例题：把七个视角连起来}
\label{ch:synthesis}

做综合题时，先问对象是什么：平面点集、二次型、线性算子，还是带算子的内积空间。再问允许的箭头是什么。这样就能判断哪些量应该保持，以及哪些反推需要额外条件。

\section{从几何图形反推代数数据}

\begin{example}{先平移，再旋转，最后求面积}
求曲线
\[
5x^2+4xy+2y^2-6x-1=0
\]
的中心、半轴长度及围成区域的面积，并写出从标准椭圆到原椭圆的全等变换。
\end{example}
\begin{solution}
把左边记作 $f(X)$，$X=(x,y)^T$。其梯度为 $2A_0X-2b$，其中 $b=(3,0)^T$。令 $A_0h=b$，解得 $h=(1,-1)^T$，因为 $A_0(1,-1)^T=(3,0)^T$。写 $X=h+Z$，展开后一次项抵消：
\[
f(h+Z)=Z^TA_0Z-h^TA_0h-1=Z^TA_0Z-4.
\]
所以中心为 $h$。再令 $Z=Q_0(u,v)^T$，标准方程为 $6u^2+v^2=4$。半轴长为 $2/\sqrt6$ 和 $2$，围成面积为 $4\pi/\sqrt6$。全等对应是 $(u,v)^T\mapsto h+Q_0(u,v)^T$；平移与正交变换都保距离。Hessian $2A_0$ 说明这是正定二次型的等值线，实分析的变量代换则给出相同的面积结论。
\end{solution}

\begin{example}{变换前后谱不同，为什么图形仍完全相同}
对比三组方程：
\[
\begin{array}{ll}
\text{甲：}&x^2+4y^2=1,\quad 2x^2+8y^2=2;\\
\text{乙：}&x^2-4y^2=0,\quad 2x^2-8y^2=0;\\
\text{丙：}&x^2+4y^2=1,\quad 2x^2+8y^2=1.
\end{array}
\]
判断点集是否全等，并说明为何不与同谱判据冲突。
\end{example}
\begin{solution}
甲组两边同除以 $2$ 后方程相同。正定椭圆若写为 $x^TAx=c$ 且 $c>0$，应把它化为 $x^T(A/c)x=1$，用于判断的是归一化矩阵 $A/c$ 的谱。甲组归一化矩阵相同。乙组的两条方程都是 $x=\pm2y$，点集也相同；但零水平不能除以 $c$ 归一化，正定单位椭圆的判据不适用。丙组两个右端都为 $1$，归一化谱不同，半长轴依次为 $1$ 和 $1/\sqrt2$，故不全等。三组对照说明：谱信息、等值常数、二次型的类型都要放在同一个命题中核对。
\end{solution}

\section{从离散谱走向连续谱}

\begin{example}{一串有特征值的算子，极限却没有特征值}
在 $L^2[0,1]$ 中，令 $s_n(t)=k/n$ 当 $k/n\le t<(k+1)/n$（$k=0,\ldots,n-1$），端点 $1$ 任意赋值。令 $M_nf=s_nf$，$Mf=tf$。求 $M_n$ 的谱，证明 $M_n$ 在算子范数下趋于 $M$，并解释极限为何可没有特征值。
\end{example}
\begin{solution}
将 $[0,1]$ 分成 $n$ 段，每段上都有非零的平方可积函数。支撑在第 $k$ 段的非零函数都是 $M_n$ 的特征函数，特征值为 $k/n$。若 $\lambda$ 不属于这 $n$ 个数，乘以 $1/(s_n-\lambda)$ 是 $(M_n-\lambda I)$ 的有界逆，因此谱恰为这 $n$ 个值。

由 $|s_n(t)-t|\le1/n$ 几乎处处，有 $\|(M_n-M)f\|_2\le \|f\|_2/n$。所以 $M_n\to M$。有限个谱点越来越密，在极限填满 $[0,1]$；对极限算子，要直接检查有界可逆性：$[0,1]$ 外可乘 $1/(t-\lambda)$ 取逆，区间内可用集中在 $\lambda$ 附近的小区间上的单位函数使 $\|(M-\lambda I)f\|_2$ 任意小。可是 $Mf=\lambda f$ 要求 $f$ 支撑在 $\{t=\lambda\}$ 上，这个集合测度为零，只能得到零函数。每个近似算子的特征空间支撑在一整段上；极限中单个取值点不再承载非零函数。
\end{solution}

\Needspace{9\baselineskip}
\section{复空间与实空间之间}

\begin{example}{一个复特征方向变成两个实方向}
令 $H=\begin{pmatrix}2&i\\-i&2\end{pmatrix}$ 作用在 $\mathbb C^2$ 上。求酉标准形，并把 $z^*Hz=1$ 看作 $\mathbb R^4$ 的全等问题。
\end{example}
\begin{solution}
$H^*=H$。分别计算可得
\[
H\binom{i}{1}=3\binom{i}{1},\qquad
H\binom{-i}{1}=\binom{-i}{1}.
\]
这两个向量在复内积下正交，各有长度 $\sqrt2$。把归一化向量组成列矩阵 $U$，有 $U^*U=I$、$U^*HU=\operatorname{diag}(3,1)$。令 $z=Uw$、$w_j=a_j+ib_j$，则
\[
z^*Hz=3(a_1^2+b_1^2)+(a_2^2+b_2^2).
\]
把每个复坐标拆成实部、虚部，复范数平方恰是四个实坐标的平方和，故 $U$ 同时给出 $\mathbb R^4$ 的一个正交变换。于是原等值集与上述实四维椭球全等，半轴长度为 $1/\sqrt3,1/\sqrt3,1,1$。一个一维复特征空间有两个实维度，这就是半轴成对出现的原因。

这里还有一层结构：复线性变换与乘 $i$ 交换；实正交变换未必如此。例如 $z\mapsto\overline z$ 在实坐标上保距，却把 $iz$ 送成 $-i\overline z$。忘掉复结构会允许更多等距箭头，这与范畴论章的忘却操作相呼应。
\end{solution}
